% Copyright (c) 2026, Ali-Eimaan. All rights reserved.
% SPDX-License-Identifier: BSD-3-Clause

% Convergence of consensus ADMM, in this repo's notation, with the assumptions that fail
% under switching topology identified explicitly.
%
% Depends on consensus_admm_derivation.tex. Compile: latexmk -pdf convergence_proof.tex
%
% Section 7 is the direct hook into the thesis proposal: the standard proof is written
% carefully, then made precise and unsparing about what it does not cover.

\documentclass[11pt]{article}
% Copyright (c) 2026, Ali-Eimaan. All rights reserved.
% SPDX-License-Identifier: BSD-3-Clause

% Shared preamble for every document in docs/derivations/.
% Included via % Copyright (c) 2026, Ali-Eimaan. All rights reserved.
% SPDX-License-Identifier: BSD-3-Clause

% Shared preamble for every document in docs/derivations/.
% Included via % Copyright (c) 2026, Ali-Eimaan. All rights reserved.
% SPDX-License-Identifier: BSD-3-Clause

% Shared preamble for every document in docs/derivations/.
% Included via \input{preamble}. Define macros here only -- never in an individual file,
% or the three documents will drift into inconsistent notation.
%
% Notation single-source-of-truth: docs/README_math.md section 1. Every macro below
% corresponds one-to-one to a row of that table; do not add a symbol there without a
% macro here (or vice versa).

\usepackage[margin=1in]{geometry}
\usepackage{amsmath, amssymb, amsthm}
\usepackage{mathtools}
\usepackage{bm}
\usepackage{algorithm}
\usepackage{algpseudocode}
\usepackage{hyperref}
\usepackage{cleveref}

% --- theorem environments ---------------------------------------------------------
\theoremstyle{plain}
\newtheorem{theorem}{Theorem}
\newtheorem{lemma}[theorem]{Lemma}
\newtheorem{proposition}[theorem]{Proposition}
\newtheorem{corollary}[theorem]{Corollary}
\theoremstyle{definition}
\newtheorem{definition}[theorem]{Definition}
\newtheorem{assumption}{Assumption}
\theoremstyle{remark}
\newtheorem{remark}[theorem]{Remark}

% --- notation ---------------------------------------------------------------------
% One macro per symbol in docs/README_math.md section 1.

\newcommand{\R}{\mathbb{R}}
\newcommand{\Nbr}{\mathcal{N}}                   % open neighborhood (subscript written as \Nbr_i)
\newcommand{\Nbrc}{\bar{\mathcal{N}}}            % closed neighborhood (subscript written as \Nbrc_i)
\newcommand{\ycopy}[2]{y_{#1}^{#2}}             % local copy
\newcommand{\zcon}[1]{z^{#1}}                   % consensus variable
\newcommand{\dual}[2]{\lambda_{#1}^{#2}}        % scaled dual (code `lam`)
\newcommand{\udual}[2]{\nu_{#1}^{#2}}           % unscaled dual = rho * scaled
\newcommand{\Lagr}{\mathcal{L}}
\newcommand{\Lrho}{\Lagr_\rho}
\newcommand{\norm}[1]{\left\lVert #1 \right\rVert}
\newcommand{\prox}{\operatorname{prox}}
\newcommand{\grad}{\nabla}
\newcommand{\diag}{\operatorname{diag}}
\newcommand{\tr}{\operatorname{tr}}
\newcommand{\Eset}{\mathcal{E}}                 % formation-edge set
\newcommand{\Gset}{\mathcal{G}}                 % communication graph
\newcommand{\Fiedler}{\lambda_2(L)}             % algebraic connectivity
\newcommand{\lmax}{\lambda_{\max}}
\newcommand{\lmin}{\lambda_{\min}}
\DeclareMathOperator*{\argmin}{arg\,min}
\DeclareMathOperator*{\argmax}{arg\,max}
. Define macros here only -- never in an individual file,
% or the three documents will drift into inconsistent notation.
%
% Notation single-source-of-truth: docs/README_math.md section 1. Every macro below
% corresponds one-to-one to a row of that table; do not add a symbol there without a
% macro here (or vice versa).

\usepackage[margin=1in]{geometry}
\usepackage{amsmath, amssymb, amsthm}
\usepackage{mathtools}
\usepackage{bm}
\usepackage{algorithm}
\usepackage{algpseudocode}
\usepackage{hyperref}
\usepackage{cleveref}

% --- theorem environments ---------------------------------------------------------
\theoremstyle{plain}
\newtheorem{theorem}{Theorem}
\newtheorem{lemma}[theorem]{Lemma}
\newtheorem{proposition}[theorem]{Proposition}
\newtheorem{corollary}[theorem]{Corollary}
\theoremstyle{definition}
\newtheorem{definition}[theorem]{Definition}
\newtheorem{assumption}{Assumption}
\theoremstyle{remark}
\newtheorem{remark}[theorem]{Remark}

% --- notation ---------------------------------------------------------------------
% One macro per symbol in docs/README_math.md section 1.

\newcommand{\R}{\mathbb{R}}
\newcommand{\Nbr}{\mathcal{N}}                   % open neighborhood (subscript written as \Nbr_i)
\newcommand{\Nbrc}{\bar{\mathcal{N}}}            % closed neighborhood (subscript written as \Nbrc_i)
\newcommand{\ycopy}[2]{y_{#1}^{#2}}             % local copy
\newcommand{\zcon}[1]{z^{#1}}                   % consensus variable
\newcommand{\dual}[2]{\lambda_{#1}^{#2}}        % scaled dual (code `lam`)
\newcommand{\udual}[2]{\nu_{#1}^{#2}}           % unscaled dual = rho * scaled
\newcommand{\Lagr}{\mathcal{L}}
\newcommand{\Lrho}{\Lagr_\rho}
\newcommand{\norm}[1]{\left\lVert #1 \right\rVert}
\newcommand{\prox}{\operatorname{prox}}
\newcommand{\grad}{\nabla}
\newcommand{\diag}{\operatorname{diag}}
\newcommand{\tr}{\operatorname{tr}}
\newcommand{\Eset}{\mathcal{E}}                 % formation-edge set
\newcommand{\Gset}{\mathcal{G}}                 % communication graph
\newcommand{\Fiedler}{\lambda_2(L)}             % algebraic connectivity
\newcommand{\lmax}{\lambda_{\max}}
\newcommand{\lmin}{\lambda_{\min}}
\DeclareMathOperator*{\argmin}{arg\,min}
\DeclareMathOperator*{\argmax}{arg\,max}
. Define macros here only -- never in an individual file,
% or the three documents will drift into inconsistent notation.
%
% Notation single-source-of-truth: docs/README_math.md section 1. Every macro below
% corresponds one-to-one to a row of that table; do not add a symbol there without a
% macro here (or vice versa).

\usepackage[margin=1in]{geometry}
\usepackage{amsmath, amssymb, amsthm}
\usepackage{mathtools}
\usepackage{bm}
\usepackage{algorithm}
\usepackage{algpseudocode}
\usepackage{hyperref}
\usepackage{cleveref}

% --- theorem environments ---------------------------------------------------------
\theoremstyle{plain}
\newtheorem{theorem}{Theorem}
\newtheorem{lemma}[theorem]{Lemma}
\newtheorem{proposition}[theorem]{Proposition}
\newtheorem{corollary}[theorem]{Corollary}
\theoremstyle{definition}
\newtheorem{definition}[theorem]{Definition}
\newtheorem{assumption}{Assumption}
\theoremstyle{remark}
\newtheorem{remark}[theorem]{Remark}

% --- notation ---------------------------------------------------------------------
% One macro per symbol in docs/README_math.md section 1.

\newcommand{\R}{\mathbb{R}}
\newcommand{\Nbr}{\mathcal{N}}                   % open neighborhood (subscript written as \Nbr_i)
\newcommand{\Nbrc}{\bar{\mathcal{N}}}            % closed neighborhood (subscript written as \Nbrc_i)
\newcommand{\ycopy}[2]{y_{#1}^{#2}}             % local copy
\newcommand{\zcon}[1]{z^{#1}}                   % consensus variable
\newcommand{\dual}[2]{\lambda_{#1}^{#2}}        % scaled dual (code `lam`)
\newcommand{\udual}[2]{\nu_{#1}^{#2}}           % unscaled dual = rho * scaled
\newcommand{\Lagr}{\mathcal{L}}
\newcommand{\Lrho}{\Lagr_\rho}
\newcommand{\norm}[1]{\left\lVert #1 \right\rVert}
\newcommand{\prox}{\operatorname{prox}}
\newcommand{\grad}{\nabla}
\newcommand{\diag}{\operatorname{diag}}
\newcommand{\tr}{\operatorname{tr}}
\newcommand{\Eset}{\mathcal{E}}                 % formation-edge set
\newcommand{\Gset}{\mathcal{G}}                 % communication graph
\newcommand{\Fiedler}{\lambda_2(L)}             % algebraic connectivity
\newcommand{\lmax}{\lambda_{\max}}
\newcommand{\lmin}{\lambda_{\min}}
\DeclareMathOperator*{\argmin}{arg\,min}
\DeclareMathOperator*{\argmax}{arg\,max}


\title{Convergence of Consensus ADMM, and Where It Fails Under Switching Topology}
\author{M. Ali Eiman}
\date{\today}

\begin{document}
\maketitle

\begin{abstract}
Classical consensus ADMM converges under a fixed connected graph and synchronous exact
solves; we state that result and its linear-rate refinement. The proof assumes a fixed
constraint set and synchronous updates, and a switching communication graph violates both.
This repository measures the consequences of those violations; characterising them is
thesis work.
\end{abstract}

\section{Assumptions}
\label{sec:assumptions}

Numbered for reference throughout the document:

\begin{assumption}[convex costs]
\label{ass:A1}
Each $f_i$ is closed, proper and convex.
\end{assumption}

\begin{assumption}[saddle point]
\label{ass:A2}
The unaugmented Lagrangian of the copy-based problem has a saddle point.
\end{assumption}

\begin{assumption}[fixed graph]
\label{ass:A3}
The communication graph is fixed and connected.
\end{assumption}

\begin{assumption}[synchrony]
\label{ass:A4}
Updates are synchronous: every agent completes iteration $k$ before any starts $k+1$.
\end{assumption}

\begin{assumption}[exact solves]
\label{ass:A5}
Every local subproblem is solved exactly.
\end{assumption}

Assumptions \ref{ass:A3}--\ref{ass:A5} are the ones this repository breaks; keep them in
mind, because Section~\ref{sec:breaks} is about exactly that.

\section{The Lyapunov function}

Fix a primal--dual optimum $(y^\star, z^\star, \lambda^\star)$ (existence is
Assumption~\ref{ass:A2}). Define
\begin{equation}
V^k = \tfrac{1}{\rho}\norm{\lambda^k - \lambda^\star}^2 + \rho\norm{z^k - z^\star}^2,
\label{eq:V}
\end{equation}
where the norms are the squared Frobenius norms over all copies. $V^k$ measures joint dual
and consensus distance from the optimum, with the $1/\rho$ and $\rho$ weights chosen so
that a single ADMM iteration drives it down by an amount proportional to the squared
residuals --- which is what the next section proves. The $x$-update does not appear in
$V^k$ directly; it enters through the residuals.

\section{Monotone decrease}

Three inequalities combine to give the descent. First, optimality of the x-update against
the optimum $y^\star$ gives
\begin{equation}
f(y^{k+1}) + \tfrac{\rho}{2}\norm{y^{k+1} - z^k + \lambda^k}^2
\le f(y^\star) + \tfrac{\rho}{2}\norm{y^\star - z^k + \lambda^k}^2.
\end{equation}
Second, because the z-update projects onto the consensus subspace,
\begin{equation}
\tfrac{\rho}{2}\norm{y^{k+1} - z^{k+1} + \lambda^k}^2
\le \tfrac{\rho}{2}\norm{y^{k+1} - z^\star + \lambda^k}^2
- \tfrac{\rho}{2}\norm{z^{k+1} - z^\star}^2.
\end{equation}
Third, the dual update $\lambda^{k+1} = \lambda^k + r^{k+1}$ with
$r^{k+1} = y^{k+1} - z^{k+1}$ gives
\begin{equation}
\tfrac{1}{\rho}\norm{\lambda^{k+1} - \lambda^\star}^2
= \tfrac{1}{\rho}\norm{\lambda^k - \lambda^\star}^2
+ \tfrac{2}{\rho}(\lambda^k - \lambda^\star)^\top r^{k+1}
+ \tfrac{1}{\rho}\norm{r^{k+1}}^2.
\end{equation}
Adding the three inequalities and using the saddle-point inequality
$f(y^{k+1}) - f(y^\star) + \lambda^{\star\top} r^{k+1} \ge 0$ yields the descent
\begin{equation}
V^{k+1} \le V^k - \rho\norm{r^{k+1}}^2 - \rho\norm{s^{k+1}}^2,
\label{eq:descent}
\end{equation}
with $s^{k+1} = z^{k+1} - z^k$ the dual residual. The $x$-update terms cancel exactly
against the saddle-point inequality; this cancellation is the technical core, and the full
derivation follows Boyd et al.~(2011, Appendix~A).

\section{Convergence of residuals and objective}

Summing \eqref{eq:descent} over $k$ gives
\begin{equation}
\rho\sum_{k=0}^{\infty}\Big(\norm{r^{k+1}}^2 + \norm{s^{k+1}}^2\Big) \le V^0 < \infty,
\end{equation}
so the residuals are square-summable and in particular $r^k \to 0$ and $s^k \to 0$. Because
$V^k$ is non-increasing and bounded below, it converges, and boundedness of the iterates
follows. With $r^k \to 0$ and $s^k \to 0$, every limit point is a saddle point, and under
Assumptions \ref{ass:A1}--\ref{ass:A5} the objective converges to the optimum. The
\emph{ergodic} (running-average) iterates satisfy
\begin{equation}
f(\bar y^K) - f^\star = O(1/K),
\qquad
\norm{\bar r^K} = O(1/K),
\end{equation}
where the bar denotes the average over the first $K$ iterates.

\section{Linear rate under strong convexity}

If each $f_i$ is additionally $\mu_f$-strongly convex with $L_f$-Lipschitz gradient, the
ergodic rate sharpens to a linear rate: there exist $c \in (0,1)$ and $M > 0$ such that
\begin{equation}
\norm{x^k - x^\star}^2 \le M\, c^{k}.
\end{equation}
The contraction $c$ depends on the graph through the extreme eigenvalues of its Laplacian
$L$: for consensus ADMM the relevant quantities are $\Fiedler = \lambda_2(L)$ and
$\lmax = \lambda_{\max}(L)$, and the optimal penalty implied by the bound scales as
\begin{equation}
\rho^\star = \Theta\Big(\sqrt{\Fiedler \cdot \lmax}\Big)
\end{equation}
in the homogeneous-strong-convexity regime. Notebook 05 sweeps $\rho$ empirically and
finds a wide, shallow basin around the value this formula predicts for the 4-agent cycle;
a stated bound with a measured comparison is worth more than the bound alone.

\section{Specialisation to this problem}

Assumptions \ref{ass:A1} and \ref{ass:A2} hold for the MPC problem. The local QPs are
convex, and strongly convex when the input weight $r > 0$ (or the velocity weight
$q_v > 0$ with full rank), because the input penalty $r\norm{U_i}^2$ is uniformly positive
definite in $U_i$ and the dynamics map $U_i \mapsto y_i^i$ is injective. Feasibility is
guaranteed by the box constraints, and a saddle point exists because the copy-based
problem is feasible and its objective is coercive on the feasible set. For the 4-agent
square example the condition number of the local Hessian is of order $10^1$--$10^2$ at the
reference tuning, which the linear-rate bound of Section 5 converts into a contraction of
roughly $0.9$--$0.99$ per iteration --- consistent with the tens of iterations observed in
notebook 05.

\section{What breaks under switching topology}
\label{sec:breaks}

Each violated assumption is treated separately; the failure modes are different in kind.

\paragraph{A3 (fixed graph).} When $\Nbrc_i$ changes, the \emph{decision variable changes
dimension}: a merge event adds copies and a split removes them. $V^k$ in \eqref{eq:V} is
defined on a space that no longer exists after the event, so monotone decrease does not
merely weaken --- it becomes \emph{ill-typed}. Relating the pre- and post-event iterates
requires a set-valued reset map relating the old variable space to the new one, and the
natural candidate --- extending the Lyapunov argument across that reset --- is exactly the
averaged/hybrid-transition-dwell object in the thesis proposal.

\paragraph{A4 (synchrony).} Under packet loss, agents update on different received $z$
vectors, so there is no single $z^k$ for $V^k$ to reference. Partially-asynchronous ADMM
results exist under a bounded-delay assumption (each agent's information is at most $\tau$
iterations stale), but they do not cover a simultaneously changing graph, and the loss
model here (independent drops, no bounded-delay guarantee) is outside even those results.

\paragraph{A5 (exact solves).} Inexact ADMM has summable-error results; the
fixed-iteration-budget inner solve used on hardware does not satisfy them. The practical
consequence is a residual floor, measured directly in notebook 04.

\paragraph{What the experiments show.} Moderate loss and moderate dwell times produce
bounded degradation, not divergence; disconnection produces structured failure --- the
split components drift apart and re-merge when reconnected. Structured failure is what
makes the hybrid treatment tractable, and this is a conjecture supported by measurement,
not a theorem. Label it as such.

\section{Open questions}

\begin{enumerate}
\item Is there a reset map across a merge/split event, and a dwell-time condition, under
  which the Lyapunov argument of Section 3 extends to the switched system?
\item Does a bounded-delay, partially-asynchronous ADMM result exist that is
  simultaneously valid under a jointly-connected switching graph?
\item What is the precise residual floor as a function of the inner iteration budget, and
  can it be certified rather than only measured?
\item Does the empirical linear rate of Section 5 hold uniformly in the graph under
  switching, or only pointwise?
\end{enumerate}

Each of these is a chapter-sized question; they are what an advisor actually reads this
document for.

% See augmented_lagrangian.tex: sources are attributed in prose, so \nocite{*} is what
% keeps the References section from rendering empty.
\nocite{*}
\bibliographystyle{plain}
\bibliography{references}
\end{document}
